%! TEX root = main.tex

\chapter{Homework 2}

\section{Equivalence to the Second-Order Wave Equation}
\label{sc:21}

The linearized shallow water equations govern the propagation of surface waves in a shallow domain \((x,t) \in (0,L) \times (0,T)\). The equivalence between the first-order coupled system and the classical second-order wave equation is established by eliminating one of the two unknowns.

\subsection{Strong Formulation}

The governing system reads:

\begin{equation}
\begin{cases}
\eta_{t}(x,t) + q_{x}(x,t) = 0 \\
q_{t}(x,t) + gH\eta_{x}(x,t) = 0
\end{cases}
\label{eq:strong_system}
\end{equation}

where \(\eta(x,t)\) is the free surface elevation, \(q(x,t)\) is the discharge, \(H\) is the mean water depth, and \(g\) is the gravitational acceleration. The system is supplemented by periodic boundary conditions:

\begin{equation}
\eta(0,t) = \eta(L,t), \quad q(0,t) = q(L,t)
\label{eq:bc_initial}
\end{equation}

\subsection{Decoupling the System}

The discharge variable \(q(x,t)\) is eliminated by cross-differentiation. Differentiating the first equation of \eqref{eq:strong_system} with respect to time \(t\):

\begin{equation}
\eta_{tt}(x,t) + q_{xt}(x,t) = 0 \quad \Longrightarrow \quad \eta_{tt}(x,t) = -q_{xt}(x,t)
\label{eq:eta_tt}
\end{equation}

Differentiating the second equation of \eqref{eq:strong_system} with respect to space \(x\):

\begin{equation}
q_{tx}(x,t) + gH\eta_{xx}(x,t) = 0 \quad \Longrightarrow \quad q_{tx}(x,t) = -gH\eta_{xx}(x,t)
\label{eq:q_tx}
\end{equation}

Assuming sufficient regularity, specifically \(q \in C^2((0,L)\times(0,T))\), Schwarz's theorem yields \(q_{xt} = q_{tx}\). Substituting \eqref{eq:q_tx} into \eqref{eq:eta_tt} results in the standard second-order wave equation:

\begin{equation}
\eta_{tt}(x,t) = gH\eta_{xx}(x,t) \quad \text{for } (x,t)\in(0,L)\times(0,T)
\label{eq:wave_eq}
\end{equation}

\subsection{Determination of Suitable Boundary Conditions}

The periodic boundary conditions on the coupled variables \((\eta, q)\) must be translated into conditions solely on \(\eta\) and its derivatives. The first condition is already explicit:

\begin{equation}
\eta(0,t) = \eta(L,t)
\label{eq:bc_eta}
\end{equation}

To convert the second condition \(q(0,t) = q(L,t)\), both sides are differentiated
with respect to \(t\), yielding \(q_{t}(0,t) = q_{t}(L,t)\). From the second equation of \eqref{eq:strong_system}, \(q_{t} = -gH\eta_{x}\). Evaluating at both boundaries and substituting:

\begin{equation}
-gH\eta_{x}(0,t) = -gH\eta_{x}(L,t) \quad \Longrightarrow \quad \eta_{x}(0,t) = \eta_{x}(L,t)
\label{eq:bc_etax}
\end{equation}

\subsection{Equivalence}

The first-order system \eqref{eq:strong_system} is therefore fully equivalent to the second-order acoustic wave equation:

\begin{equation}
\eta_{tt}(x,t) = gH\eta_{xx}(x,t), \quad (x,t)\in(0,L)\times(0,T)
\end{equation}

subject to the periodic boundary conditions \(\eta(0,t) = \eta(L,t)\) and \(\eta_{x}(0,t) = \eta_{x}(L,t)\), derived in \eqref{eq:bc_eta} and \eqref{eq:bc_etax} respectively.